\subsection{Existence and Uniqueness of Solutions of the Dirichlet Problem}\label{sec:dirichlet-existence-and-uniqueness}

\subsubsection{The Dirichlet Problem---Types of Solutions}\label{subsec:dirichlet-types-of-solutions}

We begin with a statement of the classical Dirichlet problem.

\begin{Definition}
\label{def:classical-dirichlet-solution}
Let $\Omega\subset\mathbb{R}^n$ be a bounded domain and suppose $f\in C_b(\Omega)$ is given. A function\[u\in C_b^{2k}(\Omega)\cap C_b^{2k-1}(\overline\Omega)\]is a \textbf{classical solution} of the Dirichlet problem if
\begin{equation}\label{eq:classical-dirichlet-divergence-equation}
L(\bm{x},D)u=\sum_{0\leq|\sigma|,|\gamma|\leq k}(-1)^{|\sigma|}D^\sigma(a_{\sigma\gamma}(\bm{x})D^\gamma u)=f
\end{equation}
in $\Omega$; and
\begin{equation}\label{eq:classical-dirichlet-homogeneous-boundary-conditions}
D^\alpha u=0\quad\text{for }|\alpha|\leq k-1
\end{equation}
on $\partial\Omega$.
\end{Definition}

One of the most important ideas of the modern analysis is that if you want to guarantee the existence of a solution to a problem, it is usually easier to do so in a ``bigger\textquotesingle\textquotesingle\ space of functions. This is clearly the case with the classical Dirichlet problem. Although we might expect a solution to have all of the smoothness suggested in the statement of the problem, we must relax the conditions on the solution at first so that we can use the methods of the last three chapters. The first step in relaxing the conditions on the solution is to state the problem in terms of Sobolev spaces.

\begin{Definition}
\label{def:strong-dirichlet-solution}
Let $\Omega\subset\mathbb{R}^n$ be a bounded domain and suppose $f\in L^2(\Omega)$ is given. A function\[u\in H^{2k}(\Omega)\cap H_0^k(\Omega)\]is a \textbf{strong solution} of the Dirichlet problem if
\begin{equation}\label{eq:strong-dirichlet-divergence-equation}
L(\bm{x},D)u=\sum_{0\leq|\sigma|,|\gamma|\leq k}(-1)^{|\sigma|}D^\sigma(a_{\sigma\gamma}(\bm{x})D^\gamma u)=f
\end{equation}
in $\Omega$.
\end{Definition}

Note the following.\begin{enumerate}\item We have relaxed the conditions not only on the solution $u$, but on the data $f$. The space $L^2(\Omega)$ is certainly the obvious space for $f$ once we have relaxed the conditions on $u$, so the additional generality (which includes such physically reasonable situations as discontinuous forcing functions) will come along ``for free.\textquotesingle\textquotesingle\ (In fact, we will be able to weaken the conditions on $f$ each time we relax the conditions on the solution, as we shall see below.)\item For classical solutions, the differential equation \eqref{eq:classical-dirichlet-divergence-equation} is taken to hold in a pointwise sense. For strong solutions, the differential equation \eqref{eq:strong-dirichlet-divergence-equation} is understood either in terms of equivalence classes (the right and left sides of the equation represent the same equivalence class of sequences in the $L^2(\Omega)$ norm) or in an ``almost everywhere\textquotesingle\textquotesingle\ sense (for those who have studied measure theory.)\item Instead of imposing boundary conditions \eqref{eq:classical-dirichlet-homogeneous-boundary-conditions} explicitly as we did in the classical problem, we have incorporated them into the space $H_0^k(\Omega)$ in the new problem.\item By combining the previous observations we see that the new problem is indeed a generalization of the classical problem; i.e., any classical solution of the Dirichlet problem is also a strong solution.\end{enumerate}
